\chapter{מבוא ורקע}\label{ch:intro}

\section{מה השתנה: מצילום כראיה לצילום כטענה}
במשך כמאה וחמישים שנה צילום עיתונאי שימש ראיה: אם יש תמונה, משהו כנראה קרה. זיוף צילום דרש מיומנות, זמן וציוד, ולכן היה נדיר וניתן לאיתור. מחוללי התמונות מבוססי דיפוזיה ששוחררו לציבור משנת \num{2022} ואילך ביטלו את המחסום הזה. כיום כל משתמש מקליד משפט ומקבל בתוך שניות תמונה פוטוריאליסטית של אירוע, אדם או מקום. התוצאה היא שינוי מעמד: תמונה כבר אינה ראיה, אלא טענה שצריך לאמת.

לשינוי הזה שני מחירים. הראשון הוא הפצת שקר, כשתמונה מזויפת נתפסת כאמת. השני, החמקמק יותר, הוא ״דיבידנד השקרן״: כשהציבור יודע שכל תמונה עשויה להיות מזויפת, קל לבעלי עניין לפטור גם תמונה אמיתית ומביכה בטענה ״זה בינה מלאכותית״. שני המחירים פוגעים בבסיס של העיתונות: היכולת לומר לציבור ״זה מה שקרה״ ולהיות מאמינים.

\section{מונחי יסוד}
\begin{table}[!htbp]\centering
\caption{ארבע קטגוריות של תמונות: אותו קובץ יכול להיות שקר או איור לגיטימי, לפי אופן הצגתו}\label{tab:taxonomy}
\begin{tabularx}{\linewidth}{@{}l Y l@{}}
\toprule
\textbf{קטגוריה} & \textbf{הגדרה} & \textbf{דוגמה} \\
\midrule
\textbf{א. צילום אותנטי} & צולם במצלמה; עברו עליו רק עריכות טכניות (חיתוך, איזון צבע, חשיפה) & צילום סוכנות מזירת אירוע \\
\textbf{ב. צילום שעבר מניפולציה} & צילום אמיתי שאובייקט נוסף אליו, הוסר ממנו או שונה בכלי גנרטיבי & הוספת עשן לשמיים \\
\textbf{ג. תמונה סינתטית} & נוצרה כולה במחולל, ללא צילום מקור & ״הפיצוץ בפנטגון״ \\
\textbf{ד. איור מוצהר} & תמונה סינתטית שמתויגת בבירור כאיור ואינה מציגה עצמה כתיעוד & איור לכתבת דעה \\
\bottomrule
\end{tabularx}
\end{table}

ההבחנה בין ג׳ לד׳ אינה טכנית אלא אתית: אותו קובץ בדיוק הוא שקר כשהוא מוצג כתיעוד, ולגיטימי כשהוא מוצג כאיור. לכן כל פתרון שמסתכל רק על הקובץ ולא על אופן הצגתו לקורא אינו שלם. זו ההנחה שעליה נשען השילוב בין פלטפורמה טכנולוגית לבין תקן אתי.

\section{כלכלת הקשב: למה זיוף משתלם}
אתרי חדשות ועמודי תוכן מרוויחים בעיקר מפרסום פרוגרמטי, כלומר מכל צפייה במודעה. ההכנסה תלויה בטראפיק, והטראפיק תלוי באלגוריתמים של הרשתות החברתיות ושל מנועי החיפוש. האלגוריתמים מתגמלים מעורבות: תגובות, שיתופים וזמן צפייה. תמונה רגשית, דרמטית או מפתיעה מייצרת מעורבות גבוהה יותר מתמונה מדויקת ומשעממת. נתוני \en{Meta} ממחישים את המגמה: שיעור הצפיות בפיד מפוסטים של עמודים שהמשתמש אינו עוקב אחריהם עלה מ-\num{8\%} ברבעון השני של \num{2021} ל-\num{24\%} ברבעון השלישי של \num{2023}~\cite{diresta2024}. המשמעות היא שתמונה ויראלית מגיעה לקהל חדש בלי שיש לה מוניטין להפסיד.

כך נוצרת לולאת תמריצים (איור~\ref{fig:loop}): תמונה סינתטית זולה מייצרת מעורבות, האלגוריתם מגביר אותה, הטראפיק הופך להכנסה, וההכנסה ממממנת את התמונה הבאה. לשחקן הזדוני העלות כמעט אפסית. למערכת חדשות לגיטימית הפיתוי עקיף יותר: מהירות פרסום, תמונה ״חזקה״ לכותרת, ולחץ של מתחרים שכבר פרסמו.

\begin{figure}[!htbp]
\begin{ltrdiagram}
% loop: TikZ / pgfplots source, built inside ltrdiagram (LTR)
\begin{tikzpicture}[
  st/.style={draw=acc,fill=soft,rounded corners=4pt,minimum width=3.3cm,minimum height=1.05cm,align=center,font=\sffamily\small},
  ar/.style={-{Stealth[length=7pt]},thick,acc}]
\node[st] (a) at (90:3.0)  {\he{יצירת תמונה סינתטית}\\\he{עלות כמעט אפסית}};
\node[st] (b) at (18:3.9)  {\he{מעורבות רגשית}\\\he{שיתופים ותגובות}};
\node[st] (c) at (-54:3.2) {\he{הגברה אלגוריתמית}\\\he{לקהל שאינו עוקב}};
\node[st] (d) at (-126:3.2){\he{טראפיק לאתר}\\\he{או לעמוד}};
\node[st] (e) at (162:3.9) {\he{הכנסות פרסום}\\\he{פרוגרמטי}};
\draw[ar] (a) to[bend left=12] (b);
\draw[ar] (b) to[bend left=12] (c);
\draw[ar] (c) to[bend left=12] (d);
\draw[ar] (d) to[bend left=12] (e);
\draw[ar] (e) to[bend left=12] (a);
\node[draw=acc2,fill=soft2,rounded corners=4pt,align=center,font=\sffamily\small,text width=2.8cm] at (0,0)
  {\he{נקודות שבירה ש-\en{VeriLens}}\\\he{תוקפת: אימות לפני}\\\he{פרסום וחסימת מימון}};
\end{tikzpicture}

\end{ltrdiagram}
\caption{לולאת התמריצים: כל שלב מזין את הבא, ואפשר לשבור אותה בשני מקומות, באימות ובמימון}\label{fig:loop}
\end{figure}

\section{ציר זמן: מקרי מבחן מתועדים}
איור~\ref{fig:timeline} מציג ארבעה מקרים מתועדים מהשנים \num{2023} עד \num{2024} שממחישים את סוגי הנזק. המקרים נבחרו כך שכל אחד מייצג מנגנון אחר: נזק כלכלי מיידי, גריפת טראפיק שיטתית, ניצול אסון לצורך פוליטי, והפצה המונית של תמונה סינתטית כמסר מחאה.

\begin{figure}[!htbp]
\begin{ltrdiagram}
% timeline: TikZ / pgfplots source, built inside ltrdiagram (LTR)
\begin{tikzpicture}[
  ev/.style={draw=acc,fill=white,rounded corners=3pt,text width=3.25cm,align=center,font=\sffamily\scriptsize,inner sep=4pt}]
\draw[very thick,ink,-{Stealth}] (0,0)--(15.6,0);
\foreach \x/\y/\lab/\t/\dir in {
  1.6/2023-05/{\he{פיצוץ מזויף בפנטגון}\\\he{חשבונות מאומתים הפיצו; המדדים ירדו לזמן קצר}}/a/1,
  5.6/2024-04/{\he{עמודי ספאם בפייסבוק}\\\he{125 עמודים; פוסט אחד: 40 מיליון צפיות}}/b/-1,
  9.6/2024-05/{\he{״כל העיניים על רפיח״}\\\he{תמונה סינתטית; כ-45 מיליון שיתופים}}/c/1,
  13.6/2024-10/{\he{הוריקן הלן}\\\he{״ילדה עם גור״ סינתטית שימשה למסר פוליטי}}/d/-1}{
  \fill[acc2] (\x,0) circle (3.2pt);
  \node[font=\sffamily\scriptsize\bfseries,ink] at (\x,-0.38*\dir) {\y};
  \node[ev] (\t) at (\x,1.55*\dir) {\lab};
  \draw[acc,thin] (\x,0)--(\t);
}
\end{tikzpicture}

\end{ltrdiagram}
\caption{ארבעה מקרים בשנתיים, וכל אחד מדגים סוג נזק אחר~\cite{npr2023,diresta2024,time2024,verify2024}}\label{fig:timeline}
\end{figure}

\section{למה עכשיו}
שלושה שינויים מתכנסים בשנים \num{2025} עד \num{2026} והופכים את הפרויקט מאפשרי להכרחי:
\begin{itemize}
\item \textbf{רגולציה:} חובות השקיפות בסעיף \num{50} לחוק הבינה המלאכותית האירופי חלות מ-\num{2} באוגוסט \num{2026}, כולל חובת גילוי על ״דיפ פייק״. לפי טיוטת ההנחיות של הנציבות, המבחן כולל ארבעה רכיבים ואינו דורש כוונה להטעות~\cite{euaiact2026}.
\item \textbf{תשתית:} מצלמות של \en{Leica}, \en{Sony}, \en{Canon} ו-\en{Google Pixel 10} חותמות כבר היום על תמונות בתקן \en{C2PA}~\cite{c2pa2026}. במאי \num{2026} השיקה \en{Canon} מערכת אימות למערכות חדשות שנבדקה עם \en{Reuters}~\cite{canon2026}.
\item \textbf{אמון הציבור:} רק \num{12\%} מהציבור בשש מדינות מרגישים בנוח עם חדשות שהופקו כולן בבינה מלאכותית, ו-\num{62\%} מעדיפים תוכן שהופק בידי אדם~\cite{reuters2025}. מערכת חדשות שתוכיח אמינות תזכה ביתרון תחרותי.
\end{itemize}

\begin{keybox}[מסקנה]
התשתית הטכנולוגית קיימת והחובה הרגולטורית בתוקף, אבל אין מוצר שמחבר את שתיהן לזרימת העבודה היומיומית של מערכת חדשות. את הפער הזה \en{VeriLens} סוגרת.
\end{keybox}
